\documentclass[14pt,a4paper]{extarticle}
\usepackage[left=3cm,right=2cm,top=2cm,bottom=2cm,headheight=18pt]{geometry}
\usepackage{fontspec}
\setmainfont{Liberation Serif}
\setsansfont{Liberation Serif}
\setmonofont{Liberation Mono}
\usepackage{unicode-math}
\setmathfont{texgyretermes-math.otf}
\usepackage{amsmath,graphicx,booktabs,longtable,array,enumitem}
\usepackage[normalem]{ulem}
\usepackage{titlesec,fancyhdr}
\usepackage[hidelinks,unicode]{hyperref}
\usepackage{url}
\urlstyle{same}
\setlength{\parindent}{1.27cm}
\setlength{\parskip}{0pt}
\linespread{1}
\setlength{\emergencystretch}{3em}
\clubpenalty=10000
\widowpenalty=10000
\titleformat{\section}{\normalfont\bfseries}{\thesection.}{0.5em}{}
\titlespacing*{\section}{0pt}{1.1em}{0.5em}
\setlist{nosep,leftmargin=1.27cm}
\pagestyle{fancy}
\fancyhf{}
\renewcommand{\headrulewidth}{0pt}
\renewcommand{\footrulewidth}{0pt}
\fancyhead[C]{\ifnum\value{page}>1\thepage\fi}
% Information-analytical material: O‘RQ-682 appended methodology §83 (11–14 pt), contents §84, comparison table §85.
\hyphenpenalty=10000
\exhyphenpenalty=10000
\setlength{\emergencystretch}{5em}
\renewcommand{\normalsize}{\fontsize{14bp}{16.8bp}\selectfont}
\normalsize
\newcommand{\tsize}{\fontsize{12bp}{14.4bp}\selectfont}
\newcolumntype{L}[1]{>{\raggedright\arraybackslash}p{#1}}
\setlength{\LTcapwidth}{\linewidth}
\begin{document}
\begin{flushright}
Oʻzbekiston Respublikasi Prezidenti huzuridagi\\
Ijtimoiy himoya milliy agentligi direktoriga\\[2mm]
Nusxasi: Oʻzbekiston Respublikasi\\
Iqtisodiyot va moliya vazirligiga\\[4mm]
mustaqil tadqiqotchi\\
Ashuraliyev Abduxoliq Akmal oʻgʻlidan\\
(ORCID 0009-0003-5482-5526)\\
Yashash joyi: 100017, Toshkent shahri,\\
Yunusobod tumani, Qashqar dahasi, 16-uy\\
\href{mailto:abdulxoliqashuraliyev@gmail.com}{abdulxoliqashuraliyev@gmail.com}\\[2mm]
2026-yil 1-oktabr
\end{flushright}

\begin{center}\textbf{TAKLIF}\par
\textbf{Oʻzbekiston Respublikasi Prezidentining 2026-yil 11-sentabrdagi PF-193-son Farmoni 14-bandi ijrosi yuzasidan ishlab chiqilayotgan normativ-huquqiy hujjat loyihasiga}\end{center}

Oʻzbekiston Respublikasi Konstitutsiyasining 40-moddasi, “Normativ-huquqiy hujjatlar toʻgʻrisida”gi Qonunning 20-moddasi va “Jismoniy va yuridik shaxslarning murojaatlari toʻgʻrisida”gi Qonunning 3 va 5-moddalariga muvofiq Agentlik Iqtisodiyot va moliya vazirligi bilan birgalikda 2026-yil 1-dekabrga qadar Prezident Administratsiyasiga kiritadigan, oilalarni koʻp oʻlchovli baholash asosida Ijtimoiy reyestrga kiritishning yangi tizimi toʻgʻrisidagi hujjat loyihasiga kiritish uchun takliflar kiritaman.

2027-yil 1-yanvardan boshlab oilaning toifasi balli tizim orqali avtomatik tarzda belgilanadi. Takliflar har bir oilaning baholash natijasi faqat tasdiqlangan qoidalardan kelib chiqishini, asoslari bilan saqlanishini va “Shaxsga doir maʼlumotlar toʻgʻrisida”gi Qonunning 24-moddasi talabiga muvofiq ariza beruvchiga tushuntirilishini taʼminlaydi, shuningdek Ijtimoiy reyestrni yuritish tartibi toʻgʻrisidagi amaldagi nizomdagi hisob-kitoblar bilan isbotlangan nomuvofiqliklarni bartaraf etadi. Amalga oshirish Davlat budjetidan qoʻshimcha xarajat talab qilmaydi. Oilalarning shaxsga doir maʼlumotlari soʻralmaydi.

Quyidagi har bir masala boʻyicha koʻrib chiqib, mazmunan javob berishingizni soʻrayman:
\begin{enumerate}[leftmargin=1.27cm,label=\arabic*.]
\item 1-materialning I boʻlimi qoidalarini (1–17-bandlar) yoki xuddi shu natijalarni taʼminlaydigan normalarni loyihaga kiritish.
\item Daromadi minimal isteʼmol xarajatlariga \textit{teng} boʻlgan oila Nizomning 32-bandi boʻyicha eʼtirof etilib, 39-banddagi hech qaysi toifaga kiritilmaydigan holatni 33-bandning “a” kichik bandiga oʻzgartish kiritish orqali bartaraf etish (1-material, II boʻlim, 3-band).
\item Nizomning 27-bandidagi formulada tomorqa maydonining manfiy boʻlishini istisno qilish (II boʻlim, 2-band).
\item Haqiqiy pul oʻtkazmasi oshganda hisobga olinadigan daromadning kamayishini istisno qilish va Nizomning 26-bandida solishtirish davrini koʻrsatish (II boʻlim, 1-band).
\item “Shaxsga doir maʼlumotlar toʻgʻrisida”gi Qonunning 24-moddasiga muvofiq ariza beruvchiga avtomatlashtirilgan qaror asoslari va eʼtiroz bildirish tartibini taqdim etishni taʼminlash (II boʻlim, 4-band).
\item Ogʻir hayotiy vaziyatga tushgan oilalar uchun “kambagʻallik chegarasidagi oila” toifasining oraligʻi qaysi daromadga qoʻllanilishini belgilash: har qanday talqinda Nizom 39-bandining “v” kichik bandi 35-band boʻyicha kiritiladigan oilalarning bir qismini qamrab olmaydi (2-material, 7-boʻlim).
\item Ushbu taklifni “Normativ-huquqiy hujjatlar toʻgʻrisida”gi Qonunning 23-moddasiga muvofiq loyihani tayyorlashda inobatga olish va loyiha muhokamasi materiallariga kiritish; loyiha ushbu Qonunning 24-moddasiga muvofiq normativ-huquqiy hujjatlar loyihalari muhokamasi portaliga qachon joylashtirilishini maʼlum qilish.
\end{enumerate}

“Jismoniy va yuridik shaxslarning murojaatlari toʻgʻrisida”gi Qonunning 27-moddasiga koʻra murojaatda qoʻyilgan barcha masalalar koʻrib chiqilgan va javobda har bir masala boʻyicha aniq asoslar keltirilgan taqdirda murojaat koʻrib chiqilgan hisoblanadi. Loyiha 2026-yil 1-dekabrga qadar kiritilishini inobatga olib, taklifni ushbu Qonunning 28-moddasida belgilangan bir oylik muddatda koʻrib chiqishingizni soʻrayman. Qaysi qoidalar loyihaga qaysi tahrirda kiritilganini yoki masala qaysi amaldagi norma bilan hal etilganini yoxud qoida qaysi sababga koʻra qabul qilinmaganini maʼlum qilishingizni soʻrayman.

Barcha hisob-kitoblar ilova qilingan dastur bilan qayta hisoblanadi. Nazorat misollari toʻplami va dastur bepul taqdim etiladi; ularni masʼul xodimlarga namoyish etishga tayyorman.

\noindent Ilovalar:
\begin{enumerate}[leftmargin=1.27cm,label=\arabic*.]
\item 1-material. Normativ-huquqiy hujjat loyihasiga kiritish uchun takliflar.
\item 2-material. Tushuntirish xati.
\item 3-material. Ilmiy-huquqiy asoslash va elektron ilova (nazorat hisob-kitoblari dasturi, testlar va natijalar).
\end{enumerate}
Materiallarning asosiy tahriri — rus tilida; oʻzbek va ingliz tilidagi tahrirlar ilova qilingan. 1-materialning oʻzbek tilidagi tahririda oʻzgartishlar davlat tilida bayon etilgan.

\bigskip
\noindent Ashuraliyev Abduxoliq Akmal oʻgʻli
\end{document}
